%  LaTeX support: latex@mdpi.com 
%  For support, please attach all files needed for compiling as well as the log file, and specify your operating system, LaTeX version, and LaTeX editor.

%=================================================================
\documentclass[asi,article,submit,pdftex,moreauthors]{Definitions/mdpi} 

%=================================================================

% MDPI internal commands
\firstpage{1} 
\makeatletter 
\setcounter{page}{\@firstpage} 
\makeatother
\pubvolume{1}
\issuenum{1}
\articlenumber{0}
\pubyear{2023}
\copyrightyear{2023}
\externaleditor{Academic Editor: }
\datereceived{10 May 2023} 
\daterevised{28 May 2023}
\dateaccepted{} 
\datepublished{} 
\hreflink{https://doi.org/} % If needed use \linebreak
%\doinum{}

%=================================================================
% Add packages and commands here. The following packages are loaded in our class file: fontenc, inputenc, calc, indentfirst, fancyhdr, graphicx, epstopdf, lastpage, ifthen, lineno, float, amsmath, setspace, enumitem, mathpazo, booktabs, titlesec, etoolbox, tabto, xcolor, soul, multirow, microtype, tikz, totcount, changepage, attrib, upgreek, cleveref, amsthm, hyphenat, natbib, hyperref, footmisc, url, geometry, newfloat, caption

\graphicspath{{figures/}} % path to folder with figures
\usepackage{subcaption}
\usepackage{listings}
\usepackage{xcolor}
\usepackage{gensymb} % for \textdegree
\usepackage{tabularx}
\usepackage{lipsum} 
\usepackage{makecell}
\usepackage{enumitem}

\usepackage{listings}
\usepackage{xcolor}

\definecolor{deepblue}{rgb}{0,0,0.5}
\definecolor{deepred}{rgb}{0.6,0,0}
\definecolor{deepmagenta}{RGB}{195,12,214}
\definecolor{deepgreen}{RGB}{29,122,30}
\definecolor{mintgreen}{RGB}{192,249,192}
\definecolor{codepurple}{RGB}{204, 0, 153}
\definecolor{LightCyan}{rgb}{0.88,1,1}
\definecolor{GainsboroPurple}{RGB}{247,176,247}
\definecolor{LightSlateGrey}{RGB}{124,118,118}
%    
\lstdefinestyle{mystyle}{
    commentstyle=\color{LightSlateGrey},
    keywordstyle=\color{deepgreen}\bfseries,
    emphstyle=\color{deepred},
    numberstyle=\tiny\color{deepmagenta},
    stringstyle=\color{codepurple},
    basicstyle=\ttfamily\scriptsize,
    breakatwhitespace=false,         
    breaklines=true,                 
    captionpos=b,                    
    keepspaces=true,                 
    numbers=left,                    
    numbersep=5pt,                  
    showspaces=false,                
    showstringspaces=false,
    showtabs=false,                  
    tabsize=2
}
\lstset{style=mystyle}

%=================================================================
% Full title of the paper (Capitalized)
\Title{A GRASS GIS Scripting Framework for Monitoring Changes in the Ephemeral Salt Lakes Melrhir and Merouane Chotts, Algeria}

% MDPI internal command: Title for citation in the left column
\TitleCitation{A GRASS GIS Scripting Framework for Monitoring Changes in the Ephemeral Salt Lakes Melrhir and Merouane Chotts, Algeria}

% Author Orchid ID: enter ID or remove command
\newcommand{\orcidauthorA}{0000-0002-5759-1089} % Polina Add \orcidA{} behind the author's name
\newcommand{\orcidauthorB}{0000-0002-6461-1551} % Olivier Add \orcidB{} behind the author's name

% Authors, for the paper (add full first names)
\Author{Polina Lemenkova $^{1}$*\orcidA{}}

%\longauthorlist{yes}

% MDPI internal command: Authors, for metadata in PDF
\AuthorNames{Polina Lemenkova}

% MDPI internal command: Authors, for citation in the left column
\AuthorCitation{Lemenkova, P.}
% If this is a Chicago style journal: Lastname, Firstname, Firstname Lastname, and Firstname Lastname.

% Affiliations / Addresses (Add [1] after \address if there is only one affiliation.)
\address{%
$^{1}$ \quad Laboratory of Image Synthesis and Analysis (LISA), École polytechnique de Bruxelles (Brussels Faculty of Engineering), Université Libre de Bruxelles (ULB). Building L, Campus du Solbosch, ULB -- LISA CP165/57, Avenue Franklin D. Roosevelt 50, Brussels 1050, Belgium.
}

% Contact information of the corresponding author
\corres{Correspondence: polina.lemenkova@ulb.be; Tel.: +32 471 86 04 59 (P.L.)}

% Current address and/or shared authorship
%\firstnote{Current address: Affiliation 3.} 

% Abstract (Do not insert blank lines, i.e. \\) 
\abstract{
Automated classification of satellite images is a challenging task that enables to use remote sensing data for environmental modelling of the Earth's landscapes. In this document, we implement a GRASS GIS-based framework of land cover types discrimination to identify changes in the endorheic basins of the ephemeral salt lakes Chott Melrhir and Chott Merouane, Algeria, employing embedded algorithms of image processing. This study presents a dataset of the nine Landsat 8-9 OLI/TIRS satellite images obtained from the USGS for a \textcolor{blue}{9}-year period of 2014-2022. The images were analysed for the changes of level in ephemerical lakes which experience temporal fluctuations in water level being dry for the most of the time and fed with water during rains. The unsupervised classification of the images was based on the implementation of the GRASS GIS algorithms by several modules: 'i.cluster' for generating image classes; 'i.maxlik' for classification by the maximal likelihood discriminant analysis, and auxiliary modules such as 'i.group', 'r.support', 'r.import', etc. The document includes a technical description of the scripts used for image processing with detailed comments on functionality of the GRASS GIS \textcolor{blue}{modules}. The results include the identified variations in the ephemeral salt lakes in the Algerian part of Sahara on \textcolor{blue}{the 9}-year period (2014-2022) using a time series of the Landsat OLI/TIRS multispectral images classified by the GRASS GIS. The main strength of the GRASS GIS framework is high speed, accuracy and effectiveness of the programming codes for image processing in environmental monitoring. The presented GitHub repository containing scripts used for satellite image analysis provides a reference to remote sensing data interpretation for environmental monitoring of the arid and semi-arid areas of Africa.
}

% Keywords
\keyword{Sahara Desert; Africa; salt lake; cartography; image analysis; satellite image; remote sensing; geography; GRASS GIS; mapping} 

% The fields PACS, MSC, and JEL may be left empty or commented out if not applicable
\PACS{91.10.Da; 91.10.Jf; 91.10.Sp; 91.10.Xa; 96.25.Vt; 91.10.Fc; 95.40.+s; 95.75.Qr; 95.75.Rs; 42.68.Wt}
\MSC{86A30; 86-08; 86A99; 86A04}
\JEL{ Y91; Q20; Q24; Q23; Q3; Q01; R11; O44; O13; Q5; Q51; Q55; N57; C6; C61}

%%%%%%%%%%%%%%%%%%%%%%%%%%%%%%%%%%%%%%%%%%
\begin{document}

\section{Introduction}

\subsection{Background}
Detecting changes in water level of the ephemeral salt lakes using remote sensing data is an important task in hydrological analysis of arid regions\hl{. It brings the} basis of higher level applications\hl{ }such as environmental monitoring, water resource management, climate change studies \hl{or} analysis of groundwater distribution\hl{. }Located in the northern Africa's Maghreb region in Sahara desert, salt lakes, or "chotts", are subject to temporal fluctuations in water level throughout the calendar year. Chotts remain dry for the most of the \hl{year and} are fed by water\hl{ }during winter\hl{, }occasional rains \hl{and wadi}. Monitoring\hl{ }changes in water supply \hl{of chotts }using\hl{ }is possible using satellite image classification\hl{. This} is based on \hl{evaluating }spectral reflectance of the objects visible on the\hl{ }surface \cite{Medjani}. Changes in lake contours are used as indicators for monitoring land cover changes and hydrological modelling \hl{through }variability of contours of the ephemeral stream beds.

Satellite images are a key data source in computational ecology and environmental monitoring. Various colours and brightness of the land cover classes visible on the images are determined by\hl{ }spectral reflectance of the pixels and can be used for mapping the\hl{ }Earth\hl{'s landscapes} \cite{Takorabt,Laghouag,Lemenkova202021,Lemenkova202205,Amri}. Using remote sensing data enables to detect heterogeneity in the Earth's landscapes \cite{Abbas}, record geological processes  \cite{ARAIBIA2022104455,LAMRI2016143} \hl{or} analyse environmental \hl{links} \cite{Deroin2019,data7060074} \hl{Furthermore}, spectral reflectance of the land cover types identified on the satellite image may be used to detect soil salinity \cite{abdennour2019detection}. \hl{More applications of the space borne data include} geomorphological studies and\hl{ }terrain \hl{analysis} \cite{Rebillard,Lemenkova202127,Lemenkova202105}\hl{. Finally, remote sensing data are widely used in} climate studies such as mapping evapotranspiration and droughts \cite{ijgi11090473}, and characteristics of soil layer \cite{MULDER20111}.  

This makes satellite images as a data source in many applications of Earth science such as mapping distribution of dust in arid regions \cite{Parajuli}, assessing and mapping erosion \cite{Mihi}, aquifer characterisation \cite{AOUDIA2020103920}, evaluating groundwater potential \cite{Derdour}, \hl{agricultural monitoring of crops} \cite{Benami,Butte,Daughtry,DAUGHTRY2000229}, \hl{hazard risk assessment} \cite{KUCHARCZYK2021112577}, geological investigations of rock physics, tectonics and hydrogeological simulations \cite{NEDJRAOUI2021104348}. The analysis of the satellite images can also be used for geophysical \hl{assessment} in\hl{ }desert areas \hl{of} African Sahara by integrated geological and\hl{ }spatial data modelling\hl{ }\cite{ZERROUK2017146}. 

\hl{The study goal is to create a GRASS GIS framework for image classification using scripts that yield good image processing capability, and to create a series of maps based on the automated land cover classification. }An important consideration for the use of satellite images for cartographic data analysis is that they \hl{need} to be classified into maps of land cover types for interpretation and analysis of landscapes. \hl{Manual }classification is a time-consuming process involving many iterations which might lead to errors and results in low accuracy. The problem is therefore focused on finding the effective algorithms to classification of satellite images.

Landsat 8-9 OLI/TIRS images recorded since 2013 as a continuation of the Landsat data collection form a valuable pool of big Earth geospatial data to evaluate land cover types and landscapes of the Earth. However, such datasets should be processed and classified \hl{by GIS }to enable the use of the satellite images in the environmental applications. \hl{At the same time, the use of the commercial GIS software is a subject to available license which can lead to restriction of the software usage. With this regard, the use of the freely available open source GRASS GIS for satellite image processing is advantageous. Besides, it includes the Python API and libraries as extended functionality. Therefore, this paper proposes}  using GRASS GIS \hl{as a free open source GIS for image classification. Furthermore, using} scripts \hl{of GRASS GIS enable }to automate classification of the remote sensing data. \hl{This is possible }by employing algorithms of image processing and geospatial analysis for discrimination of land cover types and detecting changes\hl{. }Besides the existing applications in the geomorphometry \cite{HOFIERKA2009387}, it can also be used for satellite image processing for environmental applications \cite{Lemenkova202125}.

\subsection{Motivation}

The need to automate classification of the satellite images \hl{raises the question} to develop tools that allow the automated image processing with reduced workload. \hl{This is aimed at building} a structured workflow for remote sensing data processing. \hl{The }existing types of classification methods \hl{such as }clustering, maximal likelihood \hl{or} K-means\hl{ }have revealed the importance of the optimisation of the algorithms of image processing. \hl{With this regards, }scripting techniques of the GRASS GIS ensure \hl{the level} of \hl{automation }and accelerates the speed of classification. \hl{Another question concerns the lack of similar studies. To the best of our knowledge, no existing studies have investigated the use of the GRASS GIS scripts for satellite image classification of the Algerian Chotts Melrhir and Merouane.}

\hl{Most} of the existing studies investigating Algerian Chotts are limited to geochemical analysis \cite{Hacini,Hacinietal2006,MOULLA201264} or geological evaluation of the saline brine environment as potential for lithium resources exploration \cite{KESLER201255,FLEXER20181188}. \hl{This is because} salt pan brines in desert regions of Sahara are known to be rich sources of lithium \cite{ZATOUT2020104566}\hl{, which raises attention in geochemical research}. Nevertheless, \hl{processing} remote sensing data for investigation of environmental changes of the water level in the saline lakes \hl{by GIS }presents a more advanced and complex approach, due to the task of image interpretation. These include the algorithms of image clustering and classification \hl{using} spectral reflectance of the pixels on the imagery.

\hl{Existing} mapping techniques \hl{used} in cartographic works on Algerian chotts and environmental monitoring \hl{use other software approaches }\cite{ABDENNOUR2020100087,OUERDACHI2012426,HADOUR2020100671}. \hl{More }examples include geological studies using traditional methods of GIS for visualisation \cite{RAHAL2021104983,LAOUINI201959,HADDANE2017228}. Most general methods of landscape monitoring in saline Saharan \hl{environment} operate by specific approaches of identification where\hl{ }surface features of arid environments are assessed based on \hl{the }geomorphologic analysis \hl{and} modelling \cite{Callot}. \hl{For instance, methods} of risk assessment of \hl{using information on concentration of }metals in water was implemented in Ramsar Site of the Chott Merouane \cite{Benhaddya}. 

The existing approaches rarely consider these methods, which makes classification of the remote sensing data \hl{by GRASS GIS }an actual process of environmental monitoring of African Sahara. In case of a satellite images, the shape and extent of land cover types for the evaluation of desertification can be\hl{ }accurately applied \cite{oussedik2003realisation}. More complex cases arise with partially eroded soil which can be also detected on the images where the mask for the erosion areas and surrounding background should be defined \cite{Toumietal}. \hl{Moreover, the} fluctuations in the lake level detected by the time series analysis \hl{on} the satellite images should include extensive data processing to derive information on land cover types. \hl{To fill in the existing gaps}, this paper contributes to this problem by presenting a GRASS GIS scripting framework for automatic classification of the satellite images. \hl{The study objective is }to identify \hl{and detect }changes in the \hl{level of saline }lake over years.

\hl{The main strength of this GRASS GIS-based framework consists in several advantages. These include the open source availability of the software which enables to freely use it in educational or research purposes, extended functionality with a wide variety of geospatial processing tools, and flexible mode Graphical User Interface (GUI) and command line mode enabling scripting. Moreover, similar to standard existing GIS, GRASS GIS supports both raster and vector processing tools, satellite image analysis and cartographic production. Such features of the GRASS GIS makes it a beneficial tool for remote sensing data analysis. For instance}, it is able to classify the land cover types \hl{based on remote sensing data }using clustering data partition and maximum-likelihood discriminant analysis classifier. \hl{Moreover}, the algorithm can identify classes over one year intervals on the \hl{satellite }images using repeated time series. Third, the workflow has a \hl{high automation} and speed of data processing\hl{, }due to scripting \hl{approach which supports machine-based data processing}.

\section{Study Area}

The study area is located in northern Algeria (Figure \ref{fig01})\hl{. A} special focus \hl{is laid} on the ephemeral salt lakes Chott Melrhir and Chott Merouane, Maghreb \cite{RAMDANI2009544,1991246,KARA2004361}. The formation of \hl{these} chotts is caused by complex geologic setting and tectonic development \hl{in} northern Algeria. \hl{This region belongs to} the North African Alpine chain \hl{which} is formed by large three mountain domains:\hl{ }Tell,\hl{ }pre-Atlas, and\hl{ }Atlas, which \hl{all }present the deformed margin of the stable African craton. The \hl{Earth's }crust \hl{in this area }undergone a complex deformation during the Alpine orogeny \cite{ASFIRANE199561}.\hl{ To the }south and east of the South Atlas Front, the Saharan topography presents \hl{a} longitudinal asymmetry and decreases in altitude eastwards. In Eastern Algeria,\hl{ }Atlas foreland corresponds to the distribution of chotts and sebkhas, which \hl{are located within the} active subsiding area with topographic depressions\hl{.}

Chotts are \hl{situated} within a large asymmetrical syncline of Northern Sahara Sedimentary Basin which was formed due to complex tectonic movements \cite{askri1995geology,Flotte}. \hl{Similar to} playas, \hl{Algerian chotts }present the endorrheic lakes located at the bottom of a large hydrographic basin \cite{Joly}. These salt lakes \hl{have} an immense area (nearly 3 M hectares) of shallow topographic depressions in northern Algeria located in the arid region of\hl{ }Sahara with a specific hydrological regime and high level of salinity \cite{chabour2018africa}. 

\begin{figure}[H]
\begin{adjustwidth}{-\extralength}{0cm}
%\centering
	\hspace{100pt}\includegraphics[width=16.0 cm]{Fig_01.jpg}
\end{adjustwidth}
\caption{Topographic map of Algeria. Software: GMT v. 6.1.1. Data sources: GEBCO and DCW. Rotated red square indicates study area with Chotts Melrhir and Merouane. Map source: author.
\label{fig01}}
\end{figure}

The topographic depression occupied by the Melrhir and Merouane Chotts (Figure \ref{fig02}) is located southwards of the Maghrebian Atlas in the Northern Sahara Sedimentary Basin. Their formation was largely influenced by vertical movements in the Atlas system during Mesozoic and Cenozoic periods, and Triassic-Lower Jurassic faults \cite{DELAMOTTE20099}. The aquifer is composed of Quaternary sediments such as sand, sandstone, clayey sand and gypsum with northwards groundwater flow orientation \cite{BOUSELSAL2022100747,BALLAIS1991209} which controls the geochemical content of groundwaters \cite{daoud1995caracterisation}.

The Algerian chotts occupy a large area mainly located in the east of the country. Their geographic situation and extent is largely influenced by a complex combination of the topographic and climatic factors: the extent of the depression\hl{ }favours the accumulation of water, sediments and dust \cite{Klose,Goudie2022}, \hl{while} a higher level of precipitation and moisture in the eastern\hl{ }Sahara\hl{ }presents the wettest region in the country \cite{Nedjraoui,Benkhaled}. \hl{Moreover}, regional environmental setting of the eastern Algeria and northern Sahara experiences the contrasting influence of the Mediterranean climate in the north and the arid Sahara in the south \cite{Ficheur,Mebarki}. 

The largest chotts of Algeria\hl{ }include the Merouane and Melrhir Chotts (Figure \ref{fig02}) \cite{HACINI2006284}. \hl{Besides multiple minor chotts of smaller size, }other existing large chotts\hl{ of Algeria }include\hl{ }Chott Ech Chergui -- the second largest chott in North Africa \cite{Maizi,habibi2013analyse,Demnati,Larnaude}, Chott El Beida (Chott El Beïdha -- Hammam Essoukhna), Chott Oum El Raneb \cite{KHAZNADAR2009369}, Chott El-Gharbi \cite{CHERCHALI2023105537} and Chott El Hodna \cite{AISSAOUI1989413}. The specific hydrologic setting of chotts consists in \hl{a highly} contrasting regime: chotts remain dry for much of the year but collect drainage water in the winter period from the occasional rains, \hl{wadi}, spring thaw waters or groundwaters \cite{Larfa,achour2014vulnerabilite}. The level of a chott surface depends on the position of the water table. \hl{Overall}, sediment deflation occurs during the decline of the water table, and sediment accumulation -- during \hl{the }water rise\hl{. }Such highly specific hydrogeological and hydrochemical setting of the Algerian chotts resulted from complex geological development of Algeria \cite{CABY1987335,BOUBEKRI2015471,DARLING2018277,BENYOUCEF2023105547}\hl{. This leads} to the formation of the saline minerals and lithium in the lake sediments \cite{HACINI20081583,MEFTAH20222105,DIAZCUEVAS2021101445}. 

\begin{figure}[H]
\makebox[1.0\linewidth][r]{%
	\begin{subfigure}[b]{.5\textwidth}
		\centering
			\includegraphics[width=7cm]{Fig_02a.jpg}
		\caption{}
	\end{subfigure}%
	\begin{subfigure}[b]{.5\textwidth}
		\centering
			\includegraphics[width=7cm]{Fig_02b.jpg}
		\caption{}
	\end{subfigure}%
}
\caption{(\textbf{a}): Melrhir Chott and Merouane Chott (also known as Felrhir Chott) enlarged on the Google Maps (2023); (\textbf{b}): Melrhir Chott and Merouane Chott on the Google Earth aerial scene (2023)}\label{fig02}
\end{figure}

The Merouane and Melrhir Chotts Algeria are located within the North-Western Sahara Aquifer System (NWSAS) \hl{-- }one of the largest transboundary aquifer systems in the world\hl{, }shared between Algeria, Tunisia, and Libya. With the size of over $1 M km^2$, the NWSAS largely contributes to the water supply of\hl{ }chotts through\hl{ }groundwaters \cite{HAKIMI2021100791}. Another source of water supply of the Melrhir and Merouane Chotts comes from \hl{the }occasional rainfalls\hl{. The} third type \hl{is} received from the spring thaw waters \hl{and wadi }from Atlas mountains. Thus, the water level in Melrhir and Merouane Chotts is supported by the three \hl{different} types of water: occasional rainfalls, spring thaw and groundwaters \cite{Ballais,fabre2005geologie,aubert1983observations}. 

The crystallized salt layer is being formed on the lake's surface during the period when\hl{ }lakes remain dry as a result of such fluctuations in the hydrology of the lake \cite{Remli}. The Algerian chotts receive water intermittently\hl{, }the lake level is subject to \hl{frequent} shrinkage and\hl{ }high evaporation rate, as salt layer accumulates in the surface of the sediment on the ground level of the lakes. The distribution of the salt differs within the salt pan with gypsum- and halite-rich deposits occupying the centre of the depressions and carbonate outliers found at \hl{the} medium level, \hl{as well as} calcareous lacustrine deposits lying on the aeolian sands \cite{GASSE19871}. \hl{Soil types is largely controlled by the }geologic structures in sedimentary basins \hl{which }include large and deep reservoirs\hl{ of Sahara contoured} by chotts\hl{ }filled by water during\hl{ }rainy periods \cite{HAMICHE2015261}. 

Specific soil properties and salinity gradients in Algerian chotts largely affect the distribution of\hl{ }flora and vegetation types around the saline lakes \cite{Neffar}. For instance, the contents of gypsum \hl{affects the }variation of \hl{the }vegetation cover of halophyte plants. Other geochemical parameters of soil in chotts such as pH, Na and Cl control the distribution of plants in moderate of low saline soils of chotts \cite{CHENCHOUNI2017660}. The biogeochemical cycles of carbon, nitrogen, sulphur and other elements within \hl{the }arid lands of chotts result\hl{ }in a specific type\hl{ }and distribution of\hl{ }species populations and microbal communities \cite{BOUTAIBA2011909,QUADRI2016119}. \hl{Further} details on geochemical setting and soil salinity \hl{are discussed in }the existing papers on the Algerian chotts\hl{ }\cite{Abdennour} \hl{with details on} hydrogeological evaluation of \hl{the }geothermal resources \cite{BENMARCE2021104285,MELOUAH2021102211,Abdelaziz}, analysis of the mineralogy of desert alluvial soils in wadi of northern Sahara \cite{Djili}, brine geochemistry and mineralogy of the Algerian playa \cite{HamdiAissa,Valles}, hydrological and environmental modelling in chotts \cite{Samie}.

\hl{Salt }pans of \hl{the }Chott Merouane and Melrhir\hl{ }present examples of sensitive wetlands in\hl{ }arid climate \hl{of Algeria}. \hl{The} agricultural and vegetation potential \hl{of chotts is largely controlled by soil in} northern Sahara \hl{which lacks} organic carbon. \hl{As a result, the vegetation of chotts is presented by a} specific type\hl{ }of flora adopted to a highly depleted dry soil of the semi-arid climate \cite{BOUBEHZIZ2020104539}. While the most areas of chotts are devoid of vegetation due to high salt concentrations in soil, there are\hl{ }occasional\hl{ }types distributed over the moist areas and edges of\hl{ }wetlands. \hl{Such} plants are mostly composed of halophytic, succulent and perennial species\hl{, e.g.,} \emph{Chenopodiaceae} family adapted to droughts and salty soils \cite{KHAZNADAR2009369}. Social \hl{advantages} from such unique landscapes include \hl{three} types of benefits: 1) salt mining in the pan; 2) palm inter cropping and palm monoculture; 3) grazing areas in a palm oasis on the border of the lake \cite{Demnati2012}. Despite the ecological importance, geochemical special setting, unique landscapes\hl{ }and\hl{ }specific biodiversity \hl{of the saline lakes, }studies on the Algerian Chotts remain scarce and relatively few. In view of this, current study presents a contribution to the investigation of the Algerian Chotts through\hl{ }image processing. 

%%%%%%%%%%%%%%%%%%%%%%%%%%%%%%%%%%%%%%%%%%
\section{Materials and Methods}

%----------- subsection ----------->
\subsection{Data}

The GRASS GIS-based image classification workflow is based on \hl{the} Landsat 8-9 dataset containing nine satellite images (Figure \ref{fig03}\hl{. The images cover} the region of Chott Melrhir and Chott Merouane, Algeria. The major metadata are listed in Table \ref{tab1}. 

\begin{table}[H]
\scriptsize
\caption{Satellite images used for computing the VIs: Landsat-8 OLI/TIRS collected from the USGS \textsuperscript{1}.\label{tab1}}
		\begin{tabularx}{\linewidth}{ p{1.6cm} | p{1.3cm} | p{6.8cm} | p{2.2cm} }
			\toprule
			\textbf{Date} & \textbf{Spacecraft} & \textbf{Landsat Product ID} & \textbf{Scene ID} \\
			\midrule
			2014/01/01 & 8 & \makecell{LC08\_L2SP\_193036\_20140101\_20200912\_02\_T1} & LC81930362014001LGN01 \\
			2015/01/04 & 8 & \makecell{LC08\_L2SP\_193036\_20150104\_20200910\_02\_T1} & LC81930362015004LGN01 \\
			2016/01/07 & 8 & \makecell{LC08\_L2SP\_193036\_20160107\_20200907\_02\_T1} & LC81930362016007LGN02 \\
			2017/01/25 & 8 & \makecell{LC08\_L2SP\_193036\_20170125\_20201013\_02\_T1} & LC81930362017025LGN02 \\
			2018/01/28 & 8 & \makecell{LC08\_L2SP\_193036\_20180128\_20200902\_02\_T1} & LC81930362018028LGN00 \\
			2019/01/15 & 8 & \makecell{LC08\_L2SP\_193036\_20190115\_20200830\_02\_T1} & LC81930362019015LGN00 \\	
			2020/01/02 & 8 & \makecell{LC08\_L2SP\_193036\_20200102\_20200824\_02\_T1} & LC81930362020002LGN00 \\
			2021/01/04 & 8 & \makecell{LC08\_L2SP\_193036\_20210104\_20210309\_02\_T1} & LC81930362021004LGN00 \\
			2022/01/15 & 9 & \makecell{LC09\_L2SP\_193036\_20220115\_20220118\_02\_T1} & LC91930362022015LGN00 \\		
			\bottomrule
		\end{tabularx}
\end{table}

The data were obtained from the USGS EarthExplorer repository always on January from 2014 to 2022. The Landsat OLI/TIRS data include spectral bands in visible and NIR ranges. \hl{The atmospheric correction with the example of code as applied as top-of-atmosphere reflectance using the GRASS GIS module  "i.landsat.toar ". It enables to calibrate the DN of the pixels on Landsat images using data on top-of-atmosphere reflectance and convert the updated image. The procedure was performed using the code in Listing} \ref{lst01}: 

\begin{lstlisting}[language=bash,caption=GRASS GIS code for atmospheric correction of the reflectance values of pixels (here: a case of the image on 2014).,style=mystyle,label={lst01}]
i.landsat.toar input=L8_2014. output=L8_2014_toar. sensor=oli8 method=dos1 date=2014-01-01 sun_elevation=28.79978934 product_date=2014-01-01 gain=HHHLHLHHL"
\end{lstlisting}

The images were downloaded \hl{as} TIFF files with \hl{ }necessary bands \hl{and} stored in \hl{a} GRASS GIS location (\hl{i.e., a }project storage folder). The images \hl{ present }a matrix of the 2D arrays of pixels assigned to \hl{the }values in a grey scale with data on spectral reflectance associated to pixels\hl{ }used for image classification. 

\begin{figure}[H]
\makebox[0.90\linewidth][r]{%
	\begin{subfigure}[b]{.40\textwidth}
		\centering
			\includegraphics[width=0.87\textwidth]{Fig_03a.jpg}
		\caption{2014}
	\end{subfigure}%
	\begin{subfigure}[b]{.40\textwidth}
		\centering
		\includegraphics[width=0.87\textwidth]{Fig_03b.jpg}
		\caption{2015}
	\end{subfigure}%
	\begin{subfigure}[b]{.40\textwidth}
		\centering
		\includegraphics[width=0.87\textwidth]{Fig_03c.jpg}
		\caption{2016}
	\end{subfigure}%
}\\
\vfill \vspace{1mm}
\makebox[0.90\linewidth][r]{%
	\begin{subfigure}[b]{.40\textwidth}
		\centering
			\includegraphics[width=0.87\textwidth]{Fig_03d.jpg}
		\caption{2017}
	\end{subfigure}%
	\begin{subfigure}[b]{.40\textwidth}
		\centering
		\includegraphics[width=0.87\textwidth]{Fig_03e.jpg}
		\caption{2018}
	\end{subfigure}%
	\begin{subfigure}[b]{.40\textwidth}
		\centering
		\includegraphics[width=0.87\textwidth]{Fig_03f.jpg}
		\caption{2019}
	\end{subfigure}%
}\\
\vfill \vspace{1mm}
\makebox[0.90\linewidth][r]{%
	\begin{subfigure}[b]{.40\textwidth}
		\centering
			\includegraphics[width=0.87\textwidth]{Fig_03g.jpg}
		\caption{2020}
	\end{subfigure}%
	\begin{subfigure}[b]{.40\textwidth}
		\centering
		\includegraphics[width=0.87\textwidth]{Fig_03h.jpg}
		\caption{2021}
	\end{subfigure}%
	\begin{subfigure}[b]{.40\textwidth}
		\centering
		\includegraphics[width=0.87\textwidth]{Fig_03i.jpg}
		\caption{2022}
	\end{subfigure}%
}
\caption{Landsat 8-9 OLI/TIRS images in natural RGB colours: lake changes for \hl{9} years (2014-2022).}\label{fig03}
\end{figure}

The series of the Landsat images demonstrated high quality of\hl{ }scenes as input data source: distinct colours, contrasting and clearly recognised areas corresponding to various land cover classes, no defects, faults or technical noise on the imagery, distinguishable landscape patterns and continuity, identical coverage of the scene over years and availability of multi-spectral bands. The images were processed, preliminary analysed and stored in a working directory for \hl{further processing by the }GRASS GIS. Additional mapping was performed using the Generic Mapping Tools (GMT) scripting toolset \cite{Wessel} following the existing technical cartographic workflow \cite{Lemenkova202206,Lemenkova202203}. The metadata on the Landsat images were collected by the EarthExplorer file with the .xml extension and summarised in Table \ref{tab1}. 

The data of the nine (9) Landsat satellite scenes used in this study include \hl{nine} consecutive years from 2014 to 2022 with scenes taken always in January on the following exact dates: \textbf{a}) 2014/01/01, (\textbf{b}) 2015/01/04, (\textbf{c}) 2016/01/07, (\textbf{d}) 2017/01/25, (\textbf{e}) 2018/01/28, (\textbf{f}) 2019/01/15, (\textbf{g}) 2020/01/02, (\textbf{h}) 2021/01/04, (\textbf{i}) 2022/01/15. Visible inspection of the images enables to see the dynamics of the salinity on the ephemeral salt lakes where bright cyan regions signify salted regions \hl{while dark blue areas indicate water, brownish and beige colour -- bare soil and sands with diverse landforms, and green -- vegetation areas}. 

%----------- subsection ----------->
\subsection{Methods}

The remote sensing data were \hl{processed} by \hl{the sequence of the }scripting methods\hl{ using several GRASS GIS modules}, Figure \ref{fig04}.  

\begin{figure}[H]
\begin{adjustwidth}{-\extralength}{0cm}
%\centering
	\hspace{150pt}\includegraphics[width=12.0 cm]{Fig_04.jpeg}
\end{adjustwidth}
\caption{Data and workflow processes used for study project. Diagram flowchart made in R mermaid library\hl{. }Source: author.\label{fig04}}
\end{figure}

The 'i.cluster' and 'i.maxlik'\hl{ }GRASS GIS modules\hl{ }enable to perform a key procedure converting raster images into \hl{the }classified objects \cite{Neteleretal2008}. The approach of \hl{the }unsupervised classification by the GRASS GIS is based on discriminating the pixels on the image based on their spectral signatures. Upon classification, the areas identified as different land cover classes are represented as geometrical primitives \hl{identified as} contoured areas that represent the extent and topological structures of the distinct land cover classes \cite{NetelerMitasova2008}. In the same way, classification enables to treat pixels as different objects according to the spectral signatures for land cover types grouping them into polygonal segments based on the defined parameters using\hl{ }clustering algorithm.

\subsection{Process flowchart}

\hl{A} general workflow consists of the following steps performed in the GRASS GIS: 1) image preprocessing by the 'i.landsat.toar' module which computes and corrects the images for the top-of-atmosphere radiance or reflectance and temperature for Landsat 8-9 OLI/TIRS; 2) Creating the group and subgroup of Landsat bands to select map layers for classification, using imagery data by the 'i.group' module;  3)\hl{ }unsupervised image classification performed by \hl{the }'i.cluster' module which creates\hl{ }cluster means and \hl{computes }covariance matrices as cluster spectral signatures; 4) image classification by the maximum-likelihood discriminant analysis classifier using the 'i.maxlik' module; 5) plotting, analysis and \hl{interpretation of} the obtained maps of the land cover types in \hl{northern} Algeria. 

Figure \ref{fig05} shows the outline of the general structure of the libraries within the GRASS GIS framework for image processing. The workflow of the GRASS GIS includes several modules that handle remote sensing data: g.region, i.group, i.cluster, i.maxlik, d.rast.leg and d.mon wx0. The GRASS GIS has several modules that actually process the data with the content and connections of these processes listed in Figure \ref{fig05}. 

\begin{figure}[H] 
	\centering
		\includegraphics[width=14cm]{Fig_05}
	\caption{The outline of the general structure of the GRASS GIS framework. Hierarchical structure of the process used for image processing with functionality of the key modules of the GRASS GIS and sequential workflow of the data processing. Flowchart made in R library 'mermaid'. Source: author.} \label{fig05}
\end{figure} 

%----------- subsection ----------->
\subsection{GRASS GIS framework workflow and process organisation}

The general workflow of the GRASS GIS algorithm for image processing is outlined stepwise as followed: 
 
\begin{enumerate}
\itemsep0em
	\item \emph{USGS -- Data } The images were collected from the USGS and organised in \hl{a} workspace of the GRASS GIS \hl{('location/mapset')}. The dataset included 9 files with each ca. 200 MB, organised for viewing and analysis of the images. 
	\begin{enumerate}
		\item \emph{EarthExplorer -- } The EarthExplorer repository was used to obtained and download the imagery, as a powerful resource for satellite \hl{images, }remote sensing data management system, and as a storage place for a large dataset of the Landsat 8-9 OLI/TIRS images.
		\item \emph{Image capture -- } Raster images were downloaded using a path to the Landsat files located in the EarthExplorer repository. The images were fetched using the access to the USGS open access files, and stored in a working folder of the GRASS GIS.		
	\end{enumerate}
	\item \emph{GMT -- Mapping } The map of the study area in the Algeria showing exact location of the Landsat images was plotted using the GEBCO/SRTM raster grids\hl{ by} the \hl{conventional }cartographic methods.	
	\item \emph{DCW -- Adding annotations} The Digital Chart of the World (DCW) Gazetteer was used to trace, mark and add annotations on the map of Algeria as parts cartographic workflow using data from the DCW Gazetteer database. The DCW data and annotation features derived based on 1:1.000.000 data were applied to Algeria with generalised annotations added on the maps by coordinates.	
	\item \emph{GRASS GIS -- Image processing} The images were processed in the GRASS GIS using a a sequence of modules: 'g.region', 'i.group', 'i.cluster', 'r.support', 'i.maxlik', 'd.mon', 'd.legend'.	
		\begin{enumerate}
			\item \emph{'i.group' -- Detection of the groups} The group and subgroups were selected by identifying visible bands in the Landsat collection and \hl{OLI/TIRS} sensor (8 or 9). This improved the time of computing and optimised the use of \hl{computer }memory \hl{(M1 chip 2020 of the MacBookAir)} during image processing.  
			\item \emph{'r.support' -- Semantic labelling --} Semantic labels were assigned using 'r.support' module for repeatability of regularly used images in time series. The labels were assigned using Landsat sensor OLI for each band correspondingly using function 'map' (e.g., map=rastername). Annotating the labels was done using 'r.support' GRASS GIS module.
			\item \emph{'i.cluster' -- Grouping} Each cluster on the image was generated by \hl{the k-means clustering} algorithm configured as a loop of \hl{codes} processing the image by the GRASS GIS. This resulted in 10 defined clusters visualised as groups of pixels on the images. This procedure was performed automatically by k-means clustering using GRASS GIS 'i.cluster' module. The following parameters were tested: input imagery group, input imagery subgroup (30 m visible Landsat bands) and signature file containing threshold levels for assignment of the pixels based on the intensity of brightness according to spectral reflectance (level of colour in pixels considered or ignored).
			\item \emph{'i.maxlik' -- Classification} The classification of the Landsat \hl{8-9} OLI/TIRS images is applied to each scene with extracted bands in visible spectrum (bands 1 to 7 with resolution 30 m). The classes on the images were adjusted using the cycles of algorithm executed iteratively until only the 10 classes of the land cover types \hl{remained}.
			\item \emph{Analysis --} Validation of errors was performed by the computed rejection probability values with pixel classification confidence levels for each map. This was necessary for the analysis of the classification \hl{accuracy for} the \hl{images covering }Chotts Melrhir and Merouane Lakes where neighbouring pixels cross each other in similar classes.
			 \item \emph{'d.mon wx0' and 'd.rast.leg' -- Visualization} The maps were visualised as several images of the land cover types based on the performed automated classification, separated by the one year gaps between each map. The sequence of maps representing the segments of land cover classes was converted into the graphical format. The maps were visualised with legends using plotting by the 'd.mon' and 'd.legend' GRASS GIS modules.	
		\end{enumerate}
	        \item \emph{Algeria -- Land cover classes} The \hl{following} 10 land cover classes were generated using features of objects on the images based on their spectral reflectance: 1) Bare areas (desert sands); 2) Salt hardpans; 3) Water bodies (Melrhir and Merouane Chotts); 4) Non-consolidated soils 5) Consolidated soils; 6) Sparse grassland; 7) Sparse shrubland; 8) Sparse trees 9) Sparse vegetation; 9) Grassland; 10) Broadleaved shrubland. 
	         \item \emph{FAO Classes -- Identifying} The features on the images and associate land cover classes were identified, annotated and compared for changes in years 2014-2022. These included visible changes in the time period for each year, and fluctuations in salt levels. The number of classes is chosen based on the Food and Agriculture Organization (FAO) Land Cover Legend Registry (LCLR), modified and adopted in such way that it well \hl{represents} the land \hl{cover types} on the \hl{images with }30-m resolution, fits easily into memory during computer processing, and is large enough to illustrate the landscape structure in the northeastern Algeria\hl{. Besides,} the vegetation parameters are chosen to ensure that\hl{ }patterns are included in the maps. 
\end{enumerate}

%----------- subsection ----------->
\subsection{Scripting in GRASS GIS}

First, a work project 'Algeria' was \hl{created and }defined using spatial location. \hl{The geospatial information on the} coordinate reference system (CRS) \hl{was derived from }the Landsat images \hl{and included the following data}: WGS84 datum and ellipsoid (EPSG:4326) UTM Projection Zone 32. The Landsat bands were \hl{then }imported into the project using 'r.import' module using the bilinear resampling method for reprojection \hl{which applies} the extent of \hl{the }current region and \hl{maintains the} resolution of \hl{the }output raster map (\hl{i.e., }30 m for Landsat bands in visible spectrum). The information on raster metadata was checked using the 'r.info' module, and the computational region was defined to match the scene using the 'g.region' module. The GRASS GIS code is shown in Listing \ref{lst02}:

\begin{lstlisting}[language=bash,caption=GRASS GIS code for importing the data (here: a case of the bands of Landsat image on 2014),style=mystyle,label={lst02}]
r.import input=/Users/polinalemenkova/grassdata/Algeria/LC08_L2SP_193036_20140101_20200912_02_T1_SR_B1.TIF output=L8_2014_01 resample=bilinear extent=region resolution=region
r.import input=/Users/polinalemenkova/grassdata/Algeria/LC08_L2SP_193036_20140101_20200912_02_T1_SR_B2.TIF output=L8_2014_02 resample=bilinear extent=region resolution=region
r.import input=/Users/polinalemenkova/grassdata/Algeria/LC08_L2SP_193036_20140101_20200912_02_T1_SR_B3.TIF output=L8_2014_03 resample=bilinear extent=region resolution=region
# repeated likewise for all the VIS bands until 7.
# check raster metadata:
r.info -r L8_2014_07
# list the imported files
g.list rast
\end{lstlisting}

The next step included \hl{the creation of} the groups and subgroups of the Landsat imagery data using GRASS GIS module 'i.group'\hl{. The groups and subgroups were then used to generate the} signature file. The aim of this step was to create raster map layers in a Landsat imagery group by assigning them to the subgroups with 30-m resolution with only necessary bands (visible spectra), \hl{following the code in }Listing \ref{lst03}:

\begin{lstlisting}[language=bash,caption=GRASS GIS code for grouping and subgrouping Landsat bands,style=mystyle,label={lst03}]
# store VIZ, NIR, MIR into group and a subgroup with 30 m resolution without TIR
i.group group=L8_2014 subgroup=res_30m \ input=L8_2014_01,L8_2014_02,L8_2014_03,L8_2014_04,L8_2014_05,L8_2014_06,L8_2014_07
\end{lstlisting}

Created group and a subgroup \hl{enabled} to run \hl{the }classification \hl{using} a combination of the raster map layers in a group for each year. \hl{The} Landsat OLI/TIRS images for several years (2014 to 2022) were inspected and processed\hl{. Grouping} and subgrouping \hl{of }the data supported \hl{the }organisation and management of data and raster map layers in a project 'Algeria'. Since current study included several images for different years,\hl{ }the semantic labels were set up for each group of bands\hl{. This was done with the aim }to re-use the signature file for the classification of different imagery groups for these years. \hl{the semantic labels were generated }before\hl{ }the signature files using 'r.support' module of the GRASS GIS\hl{, }as shown in Listing \ref{lst04}: 

\begin{lstlisting}[language=bash,caption=GRASS GIS code for creating semantic labels for the Landsat OLI/TIRS,style=mystyle,label={lst04}]
# Defining semantic labels for all Landsat OLI/TIRS bands
r.support map=L8_2014_01 semantic_label=OLI_1
r.support map=L8_2014_02 semantic_label=OLI_2
r.support map=L8_2014_03 semantic_label=OLI_3
r.support map=L8_2014_04 semantic_label=OLI_4
r.support map=L8_2014_05 semantic_label=OLI_5
r.support map=L8_2014_06 semantic_label=OLI_6
r.support map=L8_2014_07 semantic_label=OLI_7
\end{lstlisting}

The following steps of the classification \hl{included} modules 'i.cluster' and 'i.maxlik' \hl{which }used the \hl{previously defined }imagery group for each year (2014 to 2022) and subgroups of the Landsat bands with defined year and resolution. The signature file generated in the snippet of code above \hl{allowed} to classify the current imagery group only. This means\hl{ }that the data \hl{included }only\hl{ }visible Landsat bands with 30 m resolution without panchromatic and Thermal Infrared bands, which \hl{were} not necessary in this case. The methodological approach of the 'i.cluster' technique is based on reading the bands of the Landsat satellite \hl{images} and assigning pixels into \hl{the }clusters according to the difference in \hl{their }spectral reflectances. 

\hl{The} pixel clusters \hl{represent} categories \hl{(or landscape patches with 30 m size)} associated with land cover types on the Earth's shape\hl{. Therefore, }the clusters \hl{were }automatically defined by the GRASS GIS \hl{and }used for further steps in classification. Clustering approach of the GRASS GIS uses\hl{ }k-means algorithm. During the process of k-means clustering by 'i.cluster', the means of the groups where pixels are assigned are constantly \hl{changing}, as soon as new pixels are assigned and the means are recalculated iteratively to reach the maximal distances between the groups \hl{until all pixels are assigned into clusters}. The bands used for generating signatures of clusters are stored in the group $'L8\_2014'$, Listing \ref{lst05}.

\begin{lstlisting}[language=bash,caption=GRASS GIS code for creating signature file by k-means clustering algorithm,style=mystyle,label={lst05}]
# generating signature file and report for 10 classes
i.cluster group=L8_2014 subgroup=res_30m \
signaturefile=cluster_L8_2014 classes=10 reportfile=rep_clust_L8_2014.txt
\end{lstlisting}

The next step included the classification of the images by the 'i.maxlik' module\hl{. It }uses\hl{ }maximum likelihood discriminant analysis classifier as a core algorithm approach to classify the cells according to their spectral reflectances in satellite imagery. \hl{The} main approach of this algorithm is that it \hl{applies} the cluster means and covariance matrices generated during\hl{ }previous steps in \hl{a} signature file by the 'i.cluster' module. \hl{They serve} to determine the categories of spectral classes\hl{. The determination is performed }according to the \hl{degree of suitability} of each cell in the image \hl{to a given class }based on the distance to the cluster's centroids. The technical implementation was performed using Listing \ref{lst06}. 

\begin{lstlisting}[language=bash,caption=GRASS GIS code for image classification by maximal likelihood algorithm clustering,style=mystyle,label={lst06}]
i.maxlik group=L8_2014 subgroup=res_30m \
  signaturefile=cluster_L8_2014 \
  output=L8_2014_cluster_classes reject=L8_2014_cluster_reject
\end{lstlisting}

Mapping the results of the classification was done using the 'd.rast' module as shown in the following GRASS GIS code in Listing \ref{lst07}:
  
\begin{lstlisting}[language=bash,caption=GRASS GIS code for image classification by maximal likelihood algorithm clustering,style=mystyle,label={lst07}]
d.mon wx0
d.rast.leg L8_2014_cluster_classes
d.rast.leg L8_2014_cluster_reject
\end{lstlisting}

To make classification by the 'i.maxlik' more accurate, a thresholding procedure was applied by \hl{the }k-means clustering\hl{. This is implemented }using the 'i.cluster' modules aimed to separate the target image classes using the distance from the centroids. The identification of segments (objects) from \hl{the }imagery data was perform by the 'i.segment' module of GRASS GIS using Listing \ref{lst08}:

\begin{lstlisting}[language=bash,caption=GRASS GIS code for image segmentation by 'i.segment' algorithm clustering,style=mystyle,label={lst08}]
i.segment group=L8_2014 output=L8_2014_seg_14 threshold=0.4 memory=1000
\end{lstlisting}

This workflow of the GRASS GIS software\hl{ }briefly outlined above was used for classification \hl{of the }satellite images. The sequence of the utilised modules aimed to classify pixels \hl{and }sign\hl{ }them into clusters based on their features\hl{. Furthermore, this enabled to }discriminate spectral properties of the land cover objects using signature file and generate classification maps \hl{according }to \hl{the }assigned land cover classes. The landscapes categories are recognised \hl{based on the} to contrast in \hl{Digital Numbers} (DNs) \hl{of the pixels}, assigned to land cover classes and mapped using colour palettes. These solutions have advances and high level of automation which enabled to process a time series of the nine images. 

Once the images were classified in \hl{the }GRASS GIS,\hl{ }land cover classes were assigned to \hl{interprete} the variety of landscapes in \hl{a} target region of \hl{the }north-eastern Algeria. The land cover classes indicate dominant feature classes of the landscapes identified on the images that have specific characteristics and vegetation pattern unique for that given spot on the image. These include the following types derived for the FAO and adopted by generalisation for the given study area: 1) Bare areas (desert sands); 2) Salt hardpans; 3) Water bodies (Melrhir and Merouane Chotts); 4) Non-consolidated soils 5) Consolidated soils; 6) Sparse grassland; 7) Sparse shrubland; 8) Sparse trees 9) Sparse vegetation; 9) Grassland; 10) Broadleaved shrubland. The land cover classes were used as representative areas \hl{on} the classified maps to compare and distinguish\hl{ landscape dynamics }over \hl{nine} years. They also allowed to \hl{evaluate} the shrinkage of lakes and extent of \hl{the }salt areas, \hl{as well as to} analyse the \hl{extent of the patches within the }visualised landscapes. 

%%%%%%%%%%%%%%%%%%%%%%%%%%%%%%%%%%%%%%%%%%
%\pagebreak
\section{Results}

The qualitative evaluation of the water level coverage in the Melrhir and Merouane Chotts on the satellite images for each Landsat scene is based on the number of land cover classes as landscape segments, Figure \ref{fig06}.

\begin{figure}[H]
\hspace{-32pt}\makebox[1.00\linewidth][r]{%
	\begin{subfigure}[b]{.55\textwidth}
		\centering
			\includegraphics[width=7cm]{Fig_06a}
		\caption{\centering{2014}}
	\end{subfigure}%
	\begin{subfigure}[b]{.55\textwidth}
		\centering
			\includegraphics[width=7cm]{Fig_06b}
		\caption{\centering{2014}}
	\end{subfigure}%
}\\
\vfill \vspace{1mm}
\hspace{-32pt}\makebox[1.00\linewidth][r]{%
	\begin{subfigure}[b]{.55\textwidth}
		\centering
			\includegraphics[width=7cm]{Fig_06c}
		\caption{\centering{2015}}
	\end{subfigure}
	\begin{subfigure}[b]{.55\textwidth}
		\centering
			\includegraphics[width=7cm]{Fig_06d}
		\caption{\centering{2015}}
	\end{subfigure}%
}\caption{Classification the Landsat 8-9 OLI/TIRS of the Algeria with pixels classified into 10 classes and rejection probability values with pixel classification confidence levels: (\textbf{a}) 2014 map, (\textbf{b}) 2014 rejected classes (\textbf{c}) 2015 map, (\textbf{d}) 2015 rejected classes.}\label{fig06}
\end{figure}

The land cover types are automatically detected and visualised in Figure \ref{fig06} for years 2014 and 2015, Figure \ref{fig07} for years 2016 and 2017, Figure \ref{fig08} for years 2018 and 2019, Figure \ref{fig09} for years 2020 and 2021, and Figure \ref{fig10} for year 2022. The classification was succeeded for images representing nine Landsat scenes. The percentage of the pixels images with rejection probability values and pixel classification confidence levels is computed based on the reject threshold \hl{level}. The statistical data on \hl{the }classification results for processed scenes obtained from the GRASS GIS 'i.maxlik' is summarised in \hl{the} monochrome images, right to the classification maps, Figures \ref{fig06} to \ref{fig10}. 

The comparison of the classified image on 2016 (Figure \ref{fig07}, a) and 2017 (Figure \ref{fig07}, c) shows the visible increase of \hl{the }salt areas in the northern part of the study area during one calendar year, Figure \ref{fig07}. Such fluctuations show that significant changes in the hydrology reflect the sensibility of a chott as a closed hydrological system\hl{ }to regional changes in rainfall patterns, as also noted earlier \cite{Mahowald}.

\begin{figure}[H]
\hspace{-32pt}\makebox[1.00\linewidth][r]{%
	\begin{subfigure}[b]{.55\textwidth}
		\centering
			\includegraphics[width=7.3cm]{Fig_07a}
		\caption{\centering{2016}}
	\end{subfigure}%
	\begin{subfigure}[b]{.55\textwidth}
		\centering
			\includegraphics[width=7.3cm]{Fig_07b}
		\caption{\centering{2016}}
	\end{subfigure}%
}\\
\vfill \vspace{1mm}
\hspace{-32pt}\makebox[1.00\linewidth][r]{%
	\begin{subfigure}[b]{.55\textwidth}
		\centering
			\includegraphics[width=7.3cm]{Fig_07c}
		\caption{\centering{2017}}
	\end{subfigure}
	\begin{subfigure}[b]{.55\textwidth}
		\centering
			\includegraphics[width=7.3cm]{Fig_07d}
		\caption{\centering{2017}}
	\end{subfigure}%
}\caption{Classification the Landsat 8-9 OLI/TIRS of the Algeria with pixels classified into 10 classes and rejection probability values with pixel classification confidence levels: (\textbf{a}) 2016 map, (\textbf{b}) 2016 rejected classes (\textbf{c}) 2017 map, (\textbf{d}) 2017 rejected classes.}\label{fig07}
\end{figure}

The solutions evolve to halite and potassium and magnesium salts that are visible on the satellite images covering chotts for several years for comparison. For instance, when comparing the distribution of salt areas in 2018 (Figure \ref{fig08}, a) against 2019 (Figure \ref{fig08}, c), the increase in water level in the Melrhir Chott is visible in the latter \hl{image}. Besides \hl{the }salt minerals, the structure of the chotts is complicated by the occasional presence of sparse vegetation downstream the depression where the palm grove gives way to the salty soils, the sebkha. During the dry periods in summer time, the salt in the Algerian chotts is concentrated by evaporation, the pH increases \hl{along with the increase in }the level of carbonate calcium\hl{. }During wet periods or rains, a calcium carbonate precipitates, and the pH decreases while the residual alkalinity becomes negative \cite{Bourrie}. Such fluctuations affect the content of the evaporites (anhydrite) and soft sulfate minerals such as gypsum halites. 

\begin{figure}[H]
\hspace{-32pt}\makebox[1.00\linewidth][r]{%
	\begin{subfigure}[b]{.55\textwidth}
		\centering
			\includegraphics[width=7.3cm]{Fig_08a}
		\caption{\centering{2018}}
	\end{subfigure}%
	\begin{subfigure}[b]{.55\textwidth}
		\centering
			\includegraphics[width=7.3cm]{Fig_08b}
		\caption{\centering{2018}}
	\end{subfigure}%
}\\
\vfill \vspace{1mm}
\hspace{-32pt}\makebox[1.00\linewidth][r]{%
	\begin{subfigure}[b]{.55\textwidth}
		\centering
			\includegraphics[width=7.3cm]{Fig_08c}
		\caption{\centering{2019}}
	\end{subfigure}
	\begin{subfigure}[b]{.55\textwidth}
		\centering
			\includegraphics[width=7.3cm]{Fig_08d}
		\caption{\centering{2019}}
	\end{subfigure}%
}\caption{Classification the Landsat 8-9 OLI/TIRS of the Algeria with pixels classified into 10 classes and rejection probability values with pixel classification confidence levels: (\textbf{a}) 2018 map, (\textbf{b}) 2018 rejected classes (\textbf{c}) 2019 map, (\textbf{d}) 2019 rejected classes.}\label{fig08}
\end{figure}

Figure \ref{fig09} shows the comparison of the results of \hl{image }classification between\hl{ }2020 and\hl{ }2021 with the visualised classified maps and the rejected classes. Exporting files to the rejection probability values for processing\hl{ }pixel classification confidence levels has been performed using \hl{the }'i.maxlike' GRASS GIS module, which shows the probability of the correct classification of pixels for quality control. 

A higher salinity level in the Chott Melrhir in January 2020 (Figure \ref{fig09}, a) compared to January 2021 (Figure \ref{fig09} c) shows dry period and lesser atmospheric precipitation\hl{. Compared to 2020,} in\hl{ }2021 the sale lake \hl{was} fulled by water which indicated intense rainfalls. In contrast, the Merouane Chott demonstrated lesser variability in salinity which can be explained by the topographic structure of the basin and \hl{better} access to groundwaters. The automatic recognition of the salinity in both chotts and other land cover classes presented in the GRASS GIS algorithm of a maximum-likelihood discriminant analysis classifier\hl{. This }was performed using the empirical trials of pixels assignments into classes by clustering \hl{algorithm }with the most appropriate parameters of the Algerian landscapes detected for the classification of the satellite images. 

\begin{figure}[H]
\hspace{-32pt}\makebox[1.00\linewidth][r]{%
	\begin{subfigure}[b]{.55\textwidth}
		\centering
			\includegraphics[width=7.3cm]{Fig_09a}
		\caption{\centering{2020}}
	\end{subfigure}%
	\begin{subfigure}[b]{.55\textwidth}
		\centering
			\includegraphics[width=7.3cm]{Fig_09b}
		\caption{\centering{2020}}
	\end{subfigure}%
}\\
\vfill \vspace{1mm}
\hspace{-32pt}\makebox[1.00\linewidth][r]{%
	\begin{subfigure}[b]{.55\textwidth}
		\centering
			\includegraphics[width=7.3cm]{Fig_09c}
		\caption{\centering{2021}}
	\end{subfigure}
	\begin{subfigure}[b]{.55\textwidth}
		\centering
			\includegraphics[width=7.3cm]{Fig_09d}
		\caption{\centering{2021}}
	\end{subfigure}%
}\caption{Classification the Landsat 8-9 OLI/TIRS of the Algeria with pixels classified into 10 classes and rejection probability values with pixel classification confidence levels: (\textbf{a}) 2020 map, (\textbf{b}) 2020 rejected classes (\textbf{c}) 2021 map, (\textbf{d}) 2021 rejected classes.}\label{fig09}
\end{figure}

The proposed algorithm of GRASS GIS scripting \hl{presents} a sequence of several modules \hl{which }enabled to detect 10 land cover classes in the \hl{north-eastern region of the} Algeria\hl{ }by quantitative evaluations \hl{of the  landscapes. The }comparison of \hl{land cover types was based on }pixels' values\hl{ }which enabled to assigned them into cluster groups \hl{representing land cover types}. This presented a fully automated GRASS GIS-based workflow that utilises spectral reflectance values of the \hl{pixels} extracted from the satellite images for a parameter-defined analysis aimed at identification of \hl{the }land cover features, Figure \ref{fig10}. 

\begin{figure}[H]
\hspace{-32pt}\makebox[1.00\linewidth][r]{%
	\begin{subfigure}[b]{.55\textwidth}
		\centering
			\includegraphics[width=7.3cm]{Fig_10a}
		\caption{\centering{2022}}
	\end{subfigure}%
	\begin{subfigure}[b]{.55\textwidth}
		\centering
			\includegraphics[width=7.3cm]{Fig_10b}
		\caption{\centering{2022}}
	\end{subfigure}%
}\caption{Classification the Landsat 8-9 OLI/TIRS of the Algeria with pixels classified into 10 classes and rejection probability values with pixel classification confidence levels: (\textbf{a}) 2022 map, (\textbf{b}) 2022 rejected classes.}\label{fig10}
\end{figure}

%%%%%%%%%%%%%%%%%%%%%%%%%%%%%%%%%%%%%%%%%%
\section{Discussion}

\hl{In this section, we discuss the relationships between the approaches to the land cover classification by GIS. Despite the need for a standard approach to classification system, none of the current methods has been accepted as the standard, and many approaches exist. Methods of automated image classification include K-means clustering} \cite{PAOLAPATRICIA2020129}, \hl{principal component analysis} \cite{SHARMA2023100963}, \hl{hierarchical clustering} \cite{ZHANG2023129590}, \hl{segmentation} \cite{rs15082027}, \hl{object-based classification and deep learning approaches} \cite{ijgi12010014}. \hl{Among these, clustering uses algorithm that groups pixels with common characteristics into clusters representing land cover types. With this approach, the solution of coplanarity constraints has indeterminacy of geometric position of pixels with in a given land cover class. Such ambiguity can be efficiently solved by using positional information of the proposed classes that are predetermined and known based on land use data from FAO classification schemes.}

\hl{On the other hand, the second approach explicitly uses the algorithms of object oriented segmentation to solve the ambiguity of landscape patches in case of similar land cover classes classes (e.g., sparse grassland and sparse shrubland)} \cite{Lilay}. \hl{In this method, segments are arranged by grouping neighbouring pixels together based on similarity in colour and shape features. Note that the problem is the same as the deep learning process of semantic land objects detection} \cite{Yang,Kattenborn,LI2021105763}. \hl{Several approaches have been proposed to solve the problem with a single image, e.g., a composite pattern of multiple bands} \cite{SEYDI2023112961}, \hl{or as time series analysis with the assumption of dynamically developed landscape structures} \cite{YE2023113462}. \hl{Our solution is an alternative way for land cover change detection based on image classification using scripts of GRASS GIS framework}

The presented application of geospatial data processing using\hl{ }GRASS GIS for environmental monitoring of Saharan deserts and saline \hl{ephemeral }lakes was used for the two selected chotts -- the Melrhir Chott and Merouane Chott \hl{located }in Algeria. The demonstrated and explained workflow included \hl{the collection of }images from the USGS EarthExplorer repository, dataset storage, organising and analysis, and image processing by the GRASS GIS.\hl{ Creating the }spectral signatures\hl{ determined} the categories of the land cover classes to which image pixels are assigned during the classification process\hl{. The }higher-level \hl{image }interpretation process includes the identification and recognition of pixels and their assignment to these land cover classes based on the threshold defined using chi square test. \hl{This solution is based on technical application of scripts instead of using several different phases of image processing in existing methods of classification workflow which include the procedures of image preprocessing, training samples, classification, re-classification with merging and assigning classes.}

Detecting land cover classes on the images strongly depends on the contrast of \hl{the} digital numbers (DN) \hl{of pixels compared to those of }neighbour \hl{classes}. In case of clear cloud-free images and distinct landscape contours, \hl{detection} can be performed smoothly. However, \hl{high cloud coverage on the }image\hl{ or }noise can seriously affect the inspected \hl{scenes} and worsen recognition and identification of the land cover classes. Therefore, the images were selected with minimal cloudiness and other noise. The second issue concerns specific characteristics of \hl{the }Landsat 8-9 OLI/TIRS scenes \hl{which have 30-m} resolution \hl{in multi-spectral channels. This limits the }identification of the land objects\hl{ }to \hl{a spatial scale} of 30-m extent\hl{. Necessarily, }more detailed landscape patterns are ignored due to \hl{the restrictions of }resolution and image sharpness. 

In contrast to the existing state-of-the-art\hl{ }GIS (e.g., \hl{commercial }ArcGIS or Erdas Imagine), the GRASS GIS presents a free open source approach to satellite image processing. Moreover, it accurately detects variations in pixels's reflectance \hl{which enables to identify land cover classes. The assignment of pixels} into land cover classes \hl{and} discriminating signatures of the pixels \hl{is performed during clustering which computes the }covariance matrices calculated\hl{ }by the\hl{ }module 'i.cluster'. The results of the image classification performed by the GRASS GIS compared to the \hl{existing} methods present robust automated classification of \hl{the }remote sensing data such as Landsat \hl{8-9 }OLI/TIRS\hl{. Such an approach} worked accurately according to the resolution of the original image (30 m), which is applicable to all the images from the evaluated time series. 

\hl{For future applications, }the demonstrated approach of the GRASS GIS for image processing is useful \hl{for} big data \hl{processing. Possible cases include }multi-temporal interpretation of landscapes \hl{as spatially and temporally varying vegetation patterns as response to climate forcings} \cite{McDermid,Peteretal2020,Pengetal2020}. For these purposes, future applications of the presented algorithms of the GRASS GIS foresee\hl{ }the use of the presented workflow for automatic classification of other regions of the African Sahara \hl{ that experience changes over decades due to the climate effects. While this study used nine} satellite images Landsat 8-9 OLI/TIRS \hl{which }were classified using scripting methods\hl{, the extended studies may include other space borne data such as Sentinel-2A images of the MODIS temperature data}. Compared to the existing \hl{proprietary} software, the presented\hl{ }GRASS GIS-based framework \hl{is free and available for educational needs of studies and in research. The study }demonstrated high automation of image processing \hl{using scripts }which reduces the potential errors \hl{and minimised workflow routine due to} the machine-based \hl{data processing} in classification of the remote sensing data.\hl{ For replicability, the scripts used for image processing are presented in the Appendix of this paper for further reuse.}

\hl{A significant advantage of the presented method compared to the previous methods is that it requires only one script for running the process instead of selecting various separate procedures in the GIS menu. }Such effectiveness of the GRASS GIS is especially useful for multi-temporal analysis of the \hl{long-term} environmental trends\hl{ and modelling. In particular, the automated geospatial data processing is applicable for the analysis of} desertification \hl{which can be performed }using a time series of the satellite images\hl{. }Satellite image classification is central to many\hl{ }ecological applications including analysis of land cover changes, landscape monitoring, \hl{recognition of }vegetation \hl{patterns and} risk assessment\hl{. }Such applications heavily depend on accurate and automated classification of the satellite images\hl{. The} archives of Landsat\hl{ mission starting since 1970scan serve as reliable data sources for the }tasks \hl{which }require \hl{comparison} of old \hl{and new }datasets\hl{.} 

Classified images provide information for further \hl{analysis} in environmental monitoring\hl{, }and enable to acquire the geometry \hl{and extent }of \hl{the }landscapes \hl{in order }to get information \hl{in} changes over time by comparison of several images\hl{. }Therefore, classification of several Landsat OLI/TIRS images as time series is \hl{an actual} task \hl{in environmental studies}. The approach of the GRASS GIS\hl{ }demonstrated in this study\hl{ }enables to convert bitmap raster images into \hl{the }maps of land cover changes based on\hl{ }spectral reflectance of the \hl{pixels corresponding to }land cover types \hl{and} objects on the\hl{ }Earth.

%%%%%%%%%%%%%%%%%%%%%%%%%%%%%%%%%%%%%%%%%%
\section{Conclusions}
\hl{This} article \hl{presented} an advanced GRASS GIS-based algorithm developed to accurately and automatically classify bitmap images into the map of land cover types in northern Algeria based on the cell spectral reflectances in imagery data. While different software for satellite image processing exist, scripting methods \hl{of the} GRASS GIS\hl{ }remain the most accurate \hl{approach. This is} due to the high level of automation \hl{and machine-based data processing }that minimises human-induced errors and maximises the speed of \hl{the workflow}. The objective of the presented article was a classification of the satellite images Landsat OLI/TIRS covering\hl{ }recent \hl{nine years (2014-2022)} obtained from the USGS sources. The implemented methods of image analysis are presented together with the \hl{given }examples of GRASS GIS \hl{application} in image processing. 

\hl{The} identified levels of lakes accurately discriminated from the other land cover types and sand areas of the Sahara desert (Grand Erg Oriental). \hl{The }maintained geometric shape and topology of classes, automated recognition of pixel’ colours are presented by the GRASS GIS. \hl{In this way, this} work addressed the challenges of the\hl{ }satellite \hl{image classification} using\hl{ }GRASS GIS\hl{. A sequence of its several modules was used to }support\hl{ }the environmental analysis of the \hl{changes in }salt lakes \hl{of Algeria} using \hl{scripts}. The functionality of the GRASS GIS syntax was used for automated \hl{classification} of the Landsat 8-9 OLI/TIRS data \hl{as the developed workflow presented in this paper}. The \hl{modules included }'i.maxlik' \hl{used for} classification the cell spectral reflectances in imagery data, 'i.segment' to identify segment objects from imagery data and 'i.cluster' to define spectral \hl{signatures}. Auxiliary graphical modules were used \hl{additionally, for instance, }'i.landsat.toar' and others. 

The use of GRASS GIS for satellite image processing \hl{showed }high effectiveness\hl{. In} the future work, we consider developing the algorithms of the GRASS GIS to seek the specific applications in other regions of the Earth's deserts. We will also test other algorithms such as automatic detecting of classes using supervised classification \hl{and segmentation}. By comparing all the land cover classes of the satellite images obtained during the 10 years, it is possible to discriminate changes in salt lake level that illustrate \hl{ }climate changes\hl{ which affected land cover types }in the northern segment of Sahara\hl{, Algeria}. The principle of using GRASS GIS modules ' i.maxlik ' and 'i.cluster' \hl{are explained with detailed on }creating clusters and image classification\hl{. The} usage of the scripts by 'i.composite' for color composites from the Landsat bands and the example cases of the used workflow are presented in the \hl{provided }scripts. 

%%%%%%%%%%%%%%%%%%%%%%%%%%%%%%%%%%%%%%%%%%
%\authorcontributions{Supervision, conceptualization, methodology, software, resources, funding acquisition, project administration, writing---original draft preparation, methodology, software, data curation, visualization, formal analysis, validation, writing---review and editing, and investigation, P.L. }

\funding{The publication was funded by the Editorial Office of the Applied Systems Innovation, Multidisciplinary Digital Publishing Institute (MDPI), by providing 100\% discount for the APC of this manuscript.}

\institutionalreview{Not applicable.}

\informedconsent{Not applicable.}

\dataavailability{Not applicable} 

\acknowledgments{The author thanks the anonymous reviewers for reading, suggestions and comments that improved an earlier version of this manuscript.}

\conflictsofinterest{The author declares no conflict of interest.} 

\abbreviations{Abbreviations}{
The following abbreviations are used in this manuscript:\\

\noindent 
\begin{tabular}{@{}ll}
DCW & Digital Chart of the World\\
EPSG & European Petroleum Survey Group \\
FAO & Food and Agriculture Organization\\
GEBCO & General Bathymetric Chart of the Oceans \\
GIS & Geographic Information System\\
GRASS GIS & Geographic Resources Analysis Support System GIS\\
Landsat OLI/TIRS & Landsat Operational Land Imager and Thermal Infrared Sensor\\
LCLR & Land Cover Legend Registry \\
NIR & Near Infra Red\\
NWSAS & North-Western Sahara Aquifer System\\
SRTM & Shuttle Radar Topography Mission\\
USGS & United States Geological Survey\\
UTM & Universal Transverse Mercator \\
WGS84 & World Geodetic System 1984
\end{tabular}
}

%%%%%%%%%%%%%%%%%%%%%%%%%%%%%%%%%%%%%%%%%%
\begin{adjustwidth}{-\extralength}{0cm}
%\printendnotes[custom] % Un-comment to print a list of endnotes

\reftitle{References}
\nocite{*}

%=====================================
% References, variant A: external bibliography
%=====================================
\bibliography{Algeria.bib}

%%%%%%%%%%%%%%%%%%%%%%%%%%%%%%%%%%%%%%%%%%
\end{adjustwidth}
\end{document}
